\documentclass[11pt,a4paper]{article}

% -----------------------------------------------------------------------------
% Packages
% -----------------------------------------------------------------------------
\usepackage[utf8]{inputenc}
\usepackage[T1]{fontenc}
\usepackage{lmodern}
\usepackage[a4paper,margin=1in]{geometry}
\usepackage{graphicx}
\usepackage{booktabs}
\usepackage{tabularx}
\usepackage{array}
\usepackage{enumitem}
\usepackage{xcolor}
\usepackage{microtype}
\usepackage{caption}
\usepackage{float}
\usepackage{fancyhdr}
\usepackage{titlesec}
\usepackage[hidelinks]{hyperref}

% -----------------------------------------------------------------------------
% Document style
% -----------------------------------------------------------------------------
\definecolor{SEBlue}{HTML}{2F6BFF}
\definecolor{SENavy}{HTML}{16324F}
\definecolor{SETeal}{HTML}{168F7A}
\definecolor{SEPurple}{HTML}{7455D9}
\definecolor{SEOrange}{HTML}{E59A13}
\definecolor{SEGray}{HTML}{486581}
\definecolor{SELight}{HTML}{F4F7FA}

\hypersetup{
  colorlinks=true,
  linkcolor=SENavy,
  urlcolor=SEBlue,
  pdftitle={SE Agent: A Multi-Agent Coding System},
  pdfauthor={SE Agent Project}
}

\titleformat{\section}
  {\Large\bfseries\color{SENavy}}
  {\thesection}{0.7em}{}
\titleformat{\subsection}
  {\large\bfseries\color{SETeal}}
  {\thesubsection}{0.7em}{}

\pagestyle{fancy}
\fancyhf{}
\lhead{\textcolor{SEGray}{SE Agent}}
\rhead{\textcolor{SEGray}{Multi-Agent Coding System}}
\cfoot{\thepage}
\renewcommand{\headrulewidth}{0.3pt}

\setlength{\parindent}{0pt}
\setlength{\parskip}{0.65em}
\setlist[itemize]{leftmargin=1.6em,itemsep=0.3em,topsep=0.3em}
\setlist[enumerate]{leftmargin=1.8em,itemsep=0.4em,topsep=0.3em}
\renewcommand{\arraystretch}{1.25}
\graphicspath{{docs/figures/}}

\newcolumntype{L}[1]{>{\raggedright\arraybackslash}p{#1}}
\newcolumntype{Y}{>{\raggedright\arraybackslash}X}

\newcommand{\systemfigure}[3]{%
  \begin{figure}[H]
    \centering
    \IfFileExists{#1}{%
      \includegraphics[width=\textwidth]{#1}%
    }{%
      \fbox{\parbox[c][4cm][c]{0.92\textwidth}{\centering
        Figure not found. Upload \texttt{#1} with this document.}}%
    }
    \caption{#2}
    \label{#3}
  \end{figure}
}

\begin{document}

% -----------------------------------------------------------------------------
% Title
% -----------------------------------------------------------------------------
\begin{titlepage}
  \centering
  \vspace*{2.4cm}

  {\color{SEBlue}\rule{0.82\textwidth}{2pt}}\\[1.2cm]
  {\Huge\bfseries\color{SENavy} SE Agent:\\[0.25cm]
  A Multi-Agent Coding System\par}

  \vspace{0.8cm}
  {\Large\color{SEGray}
  Architecture, guardrails, sandboxing, and future direction\par}

  \vspace{1.2cm}
  {\color{SEBlue}\rule{0.82\textwidth}{2pt}}

  \vfill
  \begin{minipage}{0.76\textwidth}
    \centering
    \color{SEGray}
    SE Agent turns a software request into a planned, implemented,
    tested, and reviewed repository change. Specialized agents work
    together under one controlled workflow.
  \end{minipage}

  \vfill
  {\large\color{SENavy}Technical Documentation}
\end{titlepage}

\tableofcontents
\clearpage

% -----------------------------------------------------------------------------
\section{Objective}
% -----------------------------------------------------------------------------

SE Agent is a multi-agent software engineering assistant. Its purpose is to
complete coding tasks inside an existing repository.

The system accepts a feature request in plain language. It studies the
repository. It creates a plan. It edits the required files. It runs real
tests. It can diagnose failures and try again.

The main objectives are:

\begin{itemize}
  \item divide a large software task into clear engineering roles;
  \item make changes that follow the structure of the target repository;
  \item verify results with real commands instead of model opinion;
  \item protect user files and credentials while tools are running;
  \item provide a final report based on measured facts; and
  \item keep each run bounded, observable, and reproducible.
\end{itemize}

SE Agent is not only a chat interface. It is an execution system. Its output
is a repository change, a test result, and a clear account of what happened.

% -----------------------------------------------------------------------------
\section{Proposed Method and Architecture}
% -----------------------------------------------------------------------------

The system follows a staged workflow. Each stage has one focused role.
Information is stored in a shared, typed workflow state. Later agents use the
results of earlier agents.

The platform has two run modes:

\begin{itemize}
  \item \textbf{Thinking mode} uses the full pipeline. It includes requirements,
  repository analysis, risk analysis, architecture, planning, testing, and
  independent review.
  \item \textbf{Normal mode} uses a shorter path. It plans, implements, tests,
  and reports. It is faster, but it does not include the independent review
  stage.
\end{itemize}

\systemfigure
  {docs/figures/system-architecture.png}
  {Layered architecture of the SE Agent platform.}
  {fig:system-architecture}

\subsection{Individual Agents}

Table~\ref{tab:agents} lists every agent role in the current codebase.

\begin{table}[H]
\centering
\small
\begin{tabularx}{\textwidth}{L{0.23\textwidth}Y Y}
\toprule
\textbf{Agent} & \textbf{Main responsibility} & \textbf{Main output} \\
\midrule
Requirements Agent & Converts the request into clear needs, constraints, assumptions, and acceptance criteria. & Requirements specification \\

Repository Agent & Studies the repository structure, important files, conventions, dependencies, and test command. & Repository analysis \\

Risk Agent & Finds likely failures, security concerns, migrations, and regression risks. & Risk analysis \\

Architecture Agent & Selects the implementation approach. It defines components, interfaces, files, trade-offs, and rejected alternatives. & Architecture decision \\

Planner Agent & Converts the architecture into ordered tasks. It declares whether file changes are expected before implementation starts. & Implementation plan \\

Direct Planner Agent & Creates a practical plan directly from the request and repository. It is used in Normal mode. & Direct implementation plan \\

Coordinator Agent & Leads implementation. It receives the current request and plan. It can edit files, run commands, and ask specialists for help. & Repository changes \\

Coder Agent & Provides a reusable implementation role with coding and test tools. The current main workflow gives these tools directly to the Coordinator. & Code changes and test attempts \\

Debugger Agent & Reads failures, finds likely causes, and gives a focused repair direction. & Diagnosis \\

Security Agent & Reviews relevant code for unsafe input handling, secret exposure, path problems, and other security risks. & Security findings \\

Testing Specialist Agent & Suggests focused tests and ways to reproduce problems. It can run test commands. & Testing advice and evidence \\

Tester Agent & Runs the selected test command inside the sandbox. It does not use an LLM. & Measured test report and raw log \\

Reviewer Agent & Independently checks correctness, scope, requirements, quality, and test coverage. & Approval or requested changes \\

Report Agent & Summarizes the completed work, test status, limitations, and follow-up actions. & Final narrative report \\
\bottomrule
\end{tabularx}
\caption{Agent roles and outputs.}
\label{tab:agents}
\end{table}

\subsection{Workflow Dataflow}

The request enters through the command line or web interface. Preflight checks
run before an AI model is called. The system then captures the initial state of
the repository.

Thinking mode performs detailed analysis. Repository and risk analysis run in
parallel. Their results are joined before architecture and planning.

Normal mode creates a direct plan. Both modes then enter the same bounded
implementation loop. The Coordinator changes files. The Tester runs the real
project test command. Failed tests return to the Debugger. A successful
Thinking-mode run also passes through the Reviewer.

\systemfigure
  {docs/figures/runtime-dataflow.png}
  {Runtime dataflow for Thinking and Normal modes.}
  {fig:runtime-dataflow}

The final report is reconciled with measured state. The requested task, changed
files, test results, and iteration count come from the system. A model cannot
replace these facts with unsupported claims.

% -----------------------------------------------------------------------------
\section{Guardrails}
% -----------------------------------------------------------------------------

Guardrails limit what agents can see, change, and report. They are applied at
several layers.

\subsection{Workspace Boundary}

Every repository tool works relative to one workspace root. Paths are resolved
before use. Parent traversal and symbolic-link escapes are rejected. An agent
cannot use a normal file tool to reach outside the selected repository.

Git commands are also scoped to the workspace. This matters when the selected
workspace is a subdirectory of a larger Git repository.

\subsection{Secret and File Protection}

Direct file reads and writes reject known secret files. This includes
environment files, key files, and certificate material. Binary files are also
rejected from normal text reads.

File reads, command output, directory listings, and search results have size
limits. These limits reduce accidental context growth and denial-of-service
risks.

One limitation remains. The code search path should apply the same secret-file
check as the direct read path. This improvement is listed in
Section~\ref{sec:future}.

\subsection{Role-Based Tool Access}

Agents receive tools according to their role:

\begin{itemize}
  \item analysis and review agents mainly receive inspection tools;
  \item the Coordinator and Coder receive write and execution tools;
  \item the Tester executes one test command through the sandbox; and
  \item specialist agents are called through the Coordinator.
\end{itemize}

This design reduces unnecessary write access. The current testing tool set
still permits general commands. A stricter command allowlist would make the
boundary stronger.

\subsection{Typed Outputs}

Agent outputs use Pydantic schemas. Text length, list size, numeric range, and
allowed values are checked. Invalid structured output is rejected or retried.
This keeps the shared state predictable.

\subsection{Bounded Iteration}

The workflow stores an iteration counter and a maximum iteration limit. Every
retry passes through the iteration gate. The run ends when the limit is
reached. This prevents an endless implementation and debugging loop.

\subsection{Evidence-Based Completion}

The system records the working tree before the run. It compares the final tree
with this baseline. Pre-existing user changes are not counted as agent work.

A request can be reported as an intentional no-op only when the Planner states
in advance that no changes are expected and the tests pass. Otherwise, a green
pre-existing test suite is not enough to claim success.

Test status comes from a real process exit code. Review status comes from the
review gate. Final report facts are replaced with measured values before they
are returned to the user.

\subsection{API Guardrails}

The server can require a bearer token. Repository paths are validated before
the event stream starts. Remote callers are restricted to the configured
repository root. Loopback callers can use the local directory picker.

% -----------------------------------------------------------------------------
\section{Sandboxing}
% -----------------------------------------------------------------------------

Repository commands may execute generated or untrusted code. SE Agent therefore
runs commands through \texttt{sandbox/local.sh}. The Python sandbox wrapper
selects the strongest available isolation tier.

\begin{table}[H]
\centering
\small
\begin{tabularx}{\textwidth}{L{0.16\textwidth}Y Y}
\toprule
\textbf{Tier} & \textbf{Isolation} & \textbf{Meaning} \\
\midrule
\texttt{bwrap} & Mount, network, process, and IPC namespaces. Most of the host filesystem is hidden. The workspace is the main writable path. & Strongest supported local boundary \\

\texttt{unshare} & A user and network namespace. Network access is removed by default. The host filesystem is still visible. & Useful, but weaker filesystem protection \\

\texttt{none} & Clean environment and controlled working directory only. & Not a security boundary; it must be explicitly allowed \\
\bottomrule
\end{tabularx}
\caption{Sandbox isolation tiers.}
\label{tab:sandbox-tiers}
\end{table}

The sandbox probes the host by running each mechanism. It does not only check
whether a command exists. A mechanism that is installed but unusable is not
selected.

The main sandbox controls are:

\begin{itemize}
  \item commands start from an empty environment using \texttt{env -i};
  \item only a small environment allowlist is forwarded;
  \item model keys and other host secrets are not forwarded;
  \item each run receives a private temporary directory and temporary home;
  \item network access is disabled by default;
  \item the file-size limit is 512 MB where the host supports it;
  \item the process limit is 512 where the host supports it;
  \item wall-clock timeouts stop long-running commands;
  \item the Python wrapper kills the full process group after a timeout; and
  \item exit code 97 identifies a sandbox failure, not a failed test suite.
\end{itemize}

On Linux, \texttt{bwrap} provides the strongest protection. On Windows, the
system can enter WSL to use Linux namespaces. Without WSL, only the
\texttt{none} tier may be available. Production use should require a real
isolation tier.

% -----------------------------------------------------------------------------
\section{Technology Stack and AI Model Providers}
% -----------------------------------------------------------------------------

\subsection{Technology Stack}

\begin{table}[H]
\centering
\small
\begin{tabularx}{\textwidth}{L{0.27\textwidth}Y}
\toprule
\textbf{Technology} & \textbf{Use in SE Agent} \\
\midrule
Python 3.11 or newer & Main implementation language \\
Google Agent Development Kit 2.8 & Agent runtime, workflow graph, runner, events, tools, sessions, artifacts, and memory interfaces \\
Pydantic 2 & Typed agent outputs and workflow state validation \\
FastAPI & HTTP API, authentication checks, health endpoints, and run endpoint \\
Server-Sent Events & Live progress messages from the server to the browser \\
Uvicorn & ASGI application server \\
HTML, CSS, and JavaScript & Local browser console \\
HTTPX and Requests & Provider communication and HTTP handling \\
Git & Change inspection, repository history, and scoped diffs \\
Bash, Bubblewrap, Unshare, and WSL & Local command isolation \\
Pytest and pytest-asyncio & Automated tests \\
Ruff & Formatting-independent lint and code-quality checks \\
\bottomrule
\end{tabularx}
\caption{Main technologies used by the platform.}
\label{tab:tech-stack}
\end{table}

The default runtime services are in memory. Sessions, artifacts, and memory do
not survive a process restart. This is suitable for local development. A
persistent deployment needs external storage.

\subsection{AI Model Providers}

The provider is selected through configuration. The agent code does not need to
change when the provider changes.

\begin{table}[H]
\centering
\small
\begin{tabularx}{\textwidth}{L{0.18\textwidth}L{0.28\textwidth}Y}
\toprule
\textbf{Provider} & \textbf{Default model} & \textbf{Integration} \\
\midrule
NVIDIA NIM & \texttt{moonshotai/kimi-k3} & Uses the custom OpenAI-compatible adapter. This is the default provider. \\
Ollama & \texttt{qwen3:4b} & Uses the same OpenAI-compatible path. It supports local, keyless model execution. \\
Google Gemini & \texttt{gemini-3.6-flash} & Uses a native ADK Gemini adapter with schema cleanup, retries, model spreading, and quota-aware failover. \\
\bottomrule
\end{tabularx}
\caption{Supported AI model providers and current defaults.}
\label{tab:model-providers}
\end{table}

The OpenAI-compatible adapter also recovers structured tool calls when a weaker
model returns JSON as plain text. Provider requests have retry and error
handling. Gemini can distribute agents across a configured model list.

% -----------------------------------------------------------------------------
\section{Future Prospects and Improvements}
\label{sec:future}
% -----------------------------------------------------------------------------

The present system is a strong local prototype. The following improvements
would make it safer and more suitable for larger teams.

\begin{enumerate}
  \item \textbf{Use per-run workspaces.}
  Remove process-global workspace state. Give every run its own immutable
  repository context. This will prevent conflicts between concurrent requests.

  \item \textbf{Add persistent storage.}
  Store sessions, artifacts, test logs, and reports in a database and object
  store. Runs could then be resumed after a restart.

  \item \textbf{Strengthen secret protection.}
  Apply the secret-file policy to code search and every inspection path. Add
  secret scanning before prompts and before final output.

  \item \textbf{Use content hashes for change evidence.}
  Replace file-size and diff-count fingerprints with cryptographic content
  hashes. This will detect equal-sized rewrites reliably.

  \item \textbf{Tighten tool permissions.}
  Remove unrestricted command execution from read-only specialists. Use command
  categories, allowlists, and explicit approval for destructive operations.

  \item \textbf{Improve cross-platform isolation.}
  Add container or micro-VM execution. This would provide a consistent boundary
  on Linux, Windows, and hosted systems.

  \item \textbf{Make compaction provider-aware.}
  Use the selected provider for context summarization. The current compactor is
  tied to the OpenAI-compatible Kimi adapter.

  \item \textbf{Expand test support.}
  Parse detailed results from npm, Go, and Cargo. Add coverage measurement,
  mutation testing, and changed-test selection.

  \item \textbf{Add human approval gates.}
  Ask for approval before dependency upgrades, migrations, deletions, network
  access, or other high-impact changes.

  \item \textbf{Improve observability.}
  Record traces, token use, model cost, tool duration, retry reasons, and
  sandbox resource use. This will make production failures easier to explain.

  \item \textbf{Build a broader evaluation suite.}
  Test the platform on many languages, repository sizes, bug types, security
  tasks, and incomplete test suites. Compare quality, cost, and completion time
  across modes and providers.
\end{enumerate}

% -----------------------------------------------------------------------------
\section{Conclusion}
% -----------------------------------------------------------------------------

SE Agent separates software work into focused roles. One agent does not need to
understand, plan, implement, test, and review everything alone. The workflow
makes these responsibilities visible and controlled.

The strongest parts of the design are its real test execution, typed state,
bounded retry loops, measured final reports, and layered sandbox controls. The
two run modes also provide a clear choice between speed and assurance.

The system still needs stronger concurrency isolation, persistent services, and
stricter tool permissions before broad multi-user deployment. These changes are
practical extensions of the present architecture.

SE Agent shows how multiple AI agents can work as a small software engineering
team. It combines model reasoning with deterministic tools and verifiable
evidence. This combination is the foundation for dependable automated coding.

\end{document}
